\scanpage{53}
explication ``à priori'' de l'\deleted{construction de} \inserted{invariance} affirmée par Washnitzer-Monsky. C'est bien vrai si $S$ est de car. zéro, mais malheureusement c'est faux en car $p>0$, même si on considère pour $f$ des projections structurales de \inserted{\unsure{familles|famille} de} courbes elliptiques (= schém. elliptiques) !

Il s'agit donc pour le moment d'une relation heuristique et sentimentale, entre deux idées différentes, dont les relations exactes sont \deleted{\unsure{pas}} \inserted{loin} encore \inserted{d'avoir} été bien élucidées ......

\begin{center}
\rule{.27\linewidth}{.4pt}\\[-.6ex]
\rule{.22\linewidth}{.4pt}
\end{center}

Pour résumer, on se trouve maintenant devant un puzzle, où il s'agit de reconstituer une théorie d'ensemble \deleted{tenant compte} \deleted{et} combinant

a) Une théorie cohomologique des complexes d'opérateurs différentiels, \deleted{\unsure{qui|qui ont|pi}} (notamment en vue \inserted{d'application au complexe de DR})

b) Une étude systématique des structures stratifiées, notamment \inserted{\unsure{prises|en vue} \unsure{avec applications conceptuelles|avec explication conceptuelle|recevront des explications conceptuelles} \unsure{\deleted{satisfaisantes}|satisfaisante|topologiques}} celles qui s'introduisent

c) La théorie de Washnitzer-Monsky,
% archive sheet number: 52

